\section{Canonical status of the pre-Hugoniot core}
\label{sec:canonical-status}

This section closes the mathematical reference for the objects implemented before the
Hugoniot-family development.  It does not introduce a new wave object.  Its purpose is
to state explicitly which formulas are authoritative and which constructions are only
numerical or visual realizations of those formulas.

\subsection{Canonical chain}

For the general quadratic flux of Section~1, the finite chart of the wave manifold is
built from
\[
 r(z)=\binom z1,\qquad
 \bar U(\tau,z)=U_0(z)+\tau(-Q(z),P(z)),
\]
and
\[
 U_-=\bar U-\frac Y2r(z),\qquad
 U_+=\bar U+\frac Y2r(z).
\]
The Rankine--Hugoniot speed is the eigenvalue selected by $r(z)$.  Thus the
pre-Hugoniot core uses the following dependency chain:
\[
 \boxed{
 F\longrightarrow\mathcal W\longrightarrow
 \{\mathcal C,\mathcal E,\mathcal S^-,\mathcal S^+\}
 \longrightarrow\text{rarefaction field and its integral curves}.
 }
\]
Here $\mathcal C=\{Y=0\}$ is the characteristic surface and
$\mathcal E\subset\mathcal C$ is the coincidence locus.  The notation $\mathcal C$
and $\mathcal E$ is fixed throughout the project and must not be interchanged.

\subsection{Chart conventions}

The finite projective chart uses $(\tau,Y,z)$.  Near $z=\infty$ the reciprocal chart
uses $(\eta,\Upsilon,Z)$ with
\[
 Z=\frac1z,\qquad \eta=\frac{\tau}{Z^2},\qquad \tau=\eta Z^2.
\]
The compactified quantity
\[
 \widehat z=\frac2\pi\arctan z
\]
is a display coordinate only.  In particular, $Z=0$ is a genuine projective boundary
of the reciprocal chart and must not be identified with a fabricated finite point at
$\widehat z=\pm1$.

\subsection{Rarefaction geometry already established}

On $\mathcal C$, the lifted rarefaction field is
\[
 \dot\tau=\det(r,U_z),\qquad
 \dot z=\det(T,r),\qquad T=(-Q,P)^T,
\]
and its projection satisfies
\[
 \frac{d\mathcal R}{d\sigma}=\det(T,U_z)r(z).
\]
Consequently the state-space curve is obtained by projection of the same mathematical
orbit; no independent state-space ODE is part of the model.  Slow and fast labels are
spectral labels of orbit segments.  A regular orbit may cross $\mathcal E$ and continue
as the same orbit with the opposite slow/fast classification.

The reciprocal-chart analysis also establishes the possibility of a hyperbolic saddle
at projective infinity.  Its four local half-separatrices are represented in that chart.
Branches contained in $Z=0$ have no finite-$z$ representative.

\subsection{Three spaces}

The project distinguishes
\[
 \mathcal W\quad\text{(wave manifold)},\qquad
 \mathcal U\simeq\mathbb R^2_{(u,v)}\quad\text{(state space)},\qquad
 \mathcal X\simeq\mathbb R^2_{(x,t)}\quad\text{(physical space)}.
\]
For a point $P\in\mathcal W$, the canonical endpoint projections are
\[
 \pi_-(P)=U_-,\qquad \pi_+(P)=U_+.
\]
A physical solution is instead a map $U:\mathcal X\to\mathcal U$.  The formulas for
shock trajectories and self-similar rarefactions documented in Section~7 specify the
future bridge to physical space; they do not claim that a complete Riemann solver is
already implemented.

\subsection{Scope boundary}

The equations above are the canonical mathematical baseline for the next numerical and
three-space infrastructure increments.  Hugoniot families $\mathcal H_\pm$, composite
waves $\mathcal K_\pm$, viscous profiles, hysteresis and bifurcation objects are
\emph{not} silently defined by this closure.  Each must enter through the project
pipeline
\[
 \boxed{\text{definition}\to\text{pure implementation}\to\text{test}
 \to\text{numerics}\to\text{geometry}\to\text{visualization}.}
\]
